% ===== RESUMO EM PORTUGUÊS (300 palavras, IMRAD) =====
\chapter*{RESUMO}
\thispagestyle{empty}

% Introdução/Background (≈36 palavras, 12%)
Imagine um mundo onde cada artigo científico pudesse ser verificado — não por confiança no autor, mas por evidência rastreável. Esta dissertação torna esse mundo mais próximo. A produção científica contemporânea enfrenta uma crise de reprodutibilidade: 70\% dos pesquisadores relatam incapacidade de reproduzir experimentos publicados (BAKER, 2016). Ecossistemas multiagentes autônomos emergem como resposta, prometendo orquestrar pipelines de pesquisa com verificabilidade total.

% Método (≈99 palavras, 33%)
Esta dissertação apresenta o ecossistema OpenCode --- uma plataforma integrada de 106 habilidades (\textit{skills}), 125 agentes especializados e 41 conectores MCP (\textit{Model Context Protocol}) que implementa um pipeline completo de produção científica acadêmica. A arquitetura aplica desenvolvimento dirigido por especificação (SDD) com 282 documentos SPEC cobrindo 100\% dos componentes, desenvolvimento dirigido por testes (TDD) com 343 casos de teste em 8 suites (SPEC-025 a SPEC-032), e um motor de auto-evolução (\textit{AutoEvolve}) que documenta 23 ciclos de melhoria contínua (85$\rightarrow$99/100). O \textit{framework} metodológico inclui o protocolo TSAC (87 palavras banidas para detecção anti-IA), validação cruzada de Pearson, um orquestrador de raciocínio com 212 tipos em 27 categorias, e um ecossistema de 5 scanners epistemológicos (Noológico, Teleológico Reverso, Validação Cruzada, Convergência Polimática e Mapa de Trajetórias Evolutivas).

% Resultados (≈120 palavras, 40%)
Os resultados demonstram que o ecossistema OpenCode atinge cobertura completa de especificação (282/282 componentes), pontuação Qualis A1 de 99/100, e uma matriz de validação cruzada com 200+ conexões de afinidade entre componentes. O pipeline acadêmico MASWOS v5.0 gerou dissertações de 100+ páginas em formato ABNT/CNPq com notas de rodapé auditáveis contendo DOI, trecho original, tradução e fichamento crítico. A acurácia do classificador quântico (QML HAM10000) atingiu 89,52\%. O ciclo evolutivo mais recente (R18) implementou um sistema de economia de tokens com tripé Governança+Economia+Auditoria.

% Conclusão (≈45 palavras, 15%)
O OpenCode demonstra que ecossistemas multiagentes com SDD+\-TDD+\-Auto\-Evolve podem orquestrar pipelines científicos completos com qualidade verificável. A plataforma estabelece um novo paradigma para pesquisa reprodutível, onde cada afirmação é rastreável e cada decisão arquitetural é documentada como ADR.

\vspace{1cm}
\noindent\textbf{Palavras-chave:} Sistemas Multiagentes. Desenvolvimento Dirigido por Especificação. Desenvolvimento Dirigido por Testes. Protocolo MCP. Pipeline Acadêmico. Qualis A1. Ecossistema OpenCode. AutoEvolve.
\newpage
